\usepackage[T1]{fontenc}
\usepackage[utf8]{inputenc}
\PassOptionsToPackage{breaklinks=true}{hyperref}
\usepackage{lmodern}
\usepackage[margin=1in]{geometry}
\usepackage{amsmath}
\usepackage{amssymb}
\usepackage{amsthm}
\usepackage{mathtools}
\usepackage{graphicx}
\usepackage{tikz}
\usetikzlibrary{arrows.meta,positioning,calc,backgrounds,decorations.pathmorphing}
\newcommand{\orcidicon}{\includegraphics[width=0.32cm]{logo-orcid.pdf}}
\newcommand{\orcidA}{\href{https://orcid.org/\orcidauthorA}{\orcidicon}}
\usepackage[dvipsnames]{xcolor}
\usepackage{booktabs}
\usepackage{array}
\usepackage{tabularx}
\usepackage[font=small,labelfont=bf]{caption}
\usepackage{microtype}
\usepackage{fancyhdr}
\usepackage{hyperref}
\usepackage{cleveref}
\usepackage{xurl}
\usepackage{enumitem}

\numberwithin{equation}{section}

\hypersetup{
  colorlinks=true,
  linkcolor={blue!60!black},
  citecolor={blue!60!black},
  urlcolor={blue!60!black}
}

\raggedbottom

\newcolumntype{Y}{>{\raggedright\arraybackslash}X}
\newcolumntype{L}[1]{>{\raggedright\arraybackslash}p{#1}}
\newcommand{\code}[1]{\nolinkurl{#1}}

\setlength{\emergencystretch}{3em}
\widowpenalty=10000
\clubpenalty=10000
\displaywidowpenalty=10000

\fancypagestyle{firstpage}{%
  \fancyhf{}
  \fancyfoot[C]{\fontsize{7.5}{8}\selectfont Automorph Inc., Wilmington, DE, USA \;\texttt{ioannis@automorph.io}\; \mbox{Zenodo DOI: 10.5281/zenodo.23088991}\; \textcopyright\ 2026\; \mbox{CC-BY 4.0}}
  \renewcommand{\headrulewidth}{0pt}
  \renewcommand{\footrulewidth}{0pt}
}

\usepackage{aliascnt}
\newtheorem{theorem}{Theorem}[section]
\crefname{theorem}{Theorem}{Theorems}
\newaliascnt{proposition}{theorem}
\newtheorem{proposition}[proposition]{Proposition}
\aliascntresetthe{proposition}
\crefname{proposition}{Proposition}{Propositions}
\newaliascnt{lemma}{theorem}
\newtheorem{lemma}[lemma]{Lemma}
\aliascntresetthe{lemma}
\crefname{lemma}{Lemma}{Lemmas}
\newaliascnt{corollary}{theorem}
\newtheorem{corollary}[corollary]{Corollary}
\aliascntresetthe{corollary}
\crefname{corollary}{Corollary}{Corollarys}
\theoremstyle{definition}
\newaliascnt{definition}{theorem}
\newtheorem{definition}[definition]{Definition}
\aliascntresetthe{definition}
\crefname{definition}{Definition}{Definitions}
\newaliascnt{example}{theorem}
\newtheorem{example}[example]{Example}
\aliascntresetthe{example}
\crefname{example}{Example}{Examples}
\newaliascnt{observation}{theorem}
\newtheorem{observation}[observation]{Observation}
\aliascntresetthe{observation}
\crefname{observation}{Observation}{Observations}
\theoremstyle{remark}
\newaliascnt{remark}{theorem}
\newtheorem{remark}[remark]{Remark}
\aliascntresetthe{remark}
\crefname{remark}{Remark}{Remarks}

% Reduction-diagram styles
\tikzset{
  st/.style={inner sep=1.2pt, font=\small, fill=white, rounded corners=2pt},
  rw/.style={-{Stealth[length=5pt]}, thick},
  rws/.style={-{Stealth[length=5pt]}, thick, densely dashed},
  cellfill/.style={fill=blue!10, draw=none},
  defect/.style={draw=red!70!black, very thick, densely dotted},
  nf/.style={inner sep=1.5pt, font=\small, draw, fill=white, rounded corners=2pt},
}
\newcommand{\rto}{\to}
\newcommand{\rtos}{\mathrel{\to^{*}}}
\newcommand{\joinable}{\mathrel{\downarrow}}
